% O9. The uniform incumbent satisfies the balance condition: part (c)(iii) of the full statement of Proposition 3.
% Placement: Appendix A.4, in the proof of the full statement of Proposition 3, as the proof of part (c)(iii), after
% the proofs of parts (c)(i) and (c)(ii). The named block repeats part (c)(iii) for reading; in the appendix the
% statement appears once, inside the full statement (statements/B.tex).
% Source: referee/W1.md, equations (4) to (8) (the uniform family and the lower bound on J') and equation (10) (the
% balance condition); PLAN.md Part 5, O9; decision 4. NOTATION (c) and (d): the class mean gap is \bar\Gamma; W1's
% d, X=g_H-c_H and A=(r^2-p^2)/(2r^2) are written out.
% Definitions used: the gross profits g_j and the spread DT of Appendix A.1 (eq:a-indicators and the identity
% t(theta)+g(theta)=theta for theta>=p); gamma=g_H-g_L; J=DT(1-tau)=DT(g_H-c_H)/gamma.

\begin{named}{Proposition~\ref{prop:coexistence}(c)(iii), the uniform incumbent}
Let \(F_r=U[0,r]\) for \(r\) in a compact interval \(I\) with \(p<\min I\). Let every low-class value lie in
\([p,\min I]\), let every high-class value be at least \(\max I\), and let \(\E_H\theta\ge\max I+c_H\). Then
\(g_H(r)>g_L(r)\) at every \(r\in I\), the product \(\DT(r)\,[g_H(r)-c_H]\) is nondecreasing on \(I\), and \(J\) is
strictly increasing on \(I\).
\end{named}

\begin{proof}[Proof of Proposition~\ref{prop:coexistence}(c)(iii)]
Write \(\bar\Gamma=\E_H\theta-\E_L\theta\) for the class mean gap. Every \(r\in I\) has \(r\ge\min I>p\), so
\((r^2-p^2)/(2r^2)>0\) on \(I\). The proof has four steps.

\emph{Step 1: the auction terms.} Fix \(r\in I\) and let \(R\) be uniform on \([0,r]\). A high-class value
\(\theta\ge\max I\ge r\) wins against every incumbent value and pays \(\max\{p,R\}\), whose mean is
\(r/2+p^2/(2r)\). A low-class value \(\theta\in[p,\min I]\subseteq[p,r]\) earns
\(\E(\theta-\max\{p,R\})_+=(\theta^2-p^2)/(2r)\), and the target receives \(\theta\) minus this rent. Averaging over
each class gives
\[
g_H=\E_H\theta-\frac r2-\frac{p^2}{2r},\qquad g_L=\frac{\E_L\theta^2-p^2}{2r},\qquad
\DT=\frac r2+\frac{p^2}{2r}-\E_L\theta+\frac{\E_L\theta^2-p^2}{2r}=\frac{\E_L(r-\theta)^2}{2r},
\]
so \(\gamma=g_H-g_L=\E_H\theta-r/2-\E_L\theta^2/(2r)\) and \(\DT+\gamma=\bar\Gamma\) on \(I\). These functions are smooth
in \(r>0\), and
\[
\DT'=\frac{r^2-\E_L\theta^2}{2r^2}\ge0,\qquad g_H'=-\frac{r^2-p^2}{2r^2},\qquad \gamma'=-\DT',
\]
where \(\DT'\ge0\) because \(\E_L\theta^2\le(\min I)^2\le r^2\). Also \(\DT(r)>0\) for every \(r>\min I\), because every
low-class value is at most \(\min I\).

\emph{Step 2: three inequalities.} First, \(\E_H\theta-c_H\ge\max I\ge r\), so
\[
g_H-c_H=\E_H\theta-c_H-\frac r2-\frac{p^2}{2r}\ge\frac r2-\frac{p^2}{2r}=r\,\frac{r^2-p^2}{2r^2}>0 .
\]
Second, every low-class value lies in \([0,r]\), so
\[
r\DT'-\DT=\E_L\theta-\frac{\E_L\theta^2}{r}=\E_L\Bigl[\theta\Bigl(1-\frac\theta r\Bigr)\Bigr]\ge0 .
\]
Third, \(\E_H\theta\ge\max I+c_H\ge r+c_H\) and \(\E_L\theta^2\le r^2\), so \(\gamma\ge(r+c_H)-r/2-r^2/(2r)=c_H>0\).
Hence \(g_H>g_L\) and \(\bar\Gamma=\DT+\gamma>0\).

\emph{Step 3: the product.} By Steps 1 and 2,
\[
\bigl(\DT\,(g_H-c_H)\bigr)'=\DT'\,(g_H-c_H)-\frac{r^2-p^2}{2r^2}\,\DT\ge\frac{r^2-p^2}{2r^2}\,(r\DT'-\DT)\ge0 ,
\]
so \(\DT(r)\,[g_H(r)-c_H]\) is nondecreasing on \(I\).

\emph{Step 4: the return \(J\).} Since \(J=\DT\,(g_H-c_H)/\gamma\) and \(\gamma'=-\DT'\),
\begin{align*}
J'&=\frac{\bigl(\DT\,(g_H-c_H)\bigr)'\gamma+\DT\,(g_H-c_H)\,\DT'}{\gamma^2}
=\frac{\bar\Gamma\,(g_H-c_H)\,\DT'-\frac{r^2-p^2}{2r^2}\,\DT\,\gamma}{\gamma^2}\\
&\ge\frac{r^2-p^2}{2r^2}\cdot\frac{\bar\Gamma\,r\DT'-\DT\,\gamma}{\gamma^2}
\ge\frac{r^2-p^2}{2r^2}\cdot\frac{\DT\,(\bar\Gamma-\gamma)}{\gamma^2}=\frac{(r^2-p^2)\,\DT^2}{2r^2\gamma^2}.
\end{align*}
The first inequality uses \(g_H-c_H\ge r\,(r^2-p^2)/(2r^2)\) and \(\DT'\ge0\); the second uses \(r\DT'\ge\DT\) and
\(\bar\Gamma>0\). The bound is nonnegative on \(I\) and positive wherever \(\DT>0\), which includes every \(r>\min I\) by
Step 1, so \(J\) is strictly increasing on \(I\).
\end{proof}

For point-mass classes at \(h\) and \(\ell\), the hypotheses of part (c)(iii) read \(\ell\in[p,\min I]\),
\(p<\min I\) and \(\max I+c_H\le h\).
